% !TEX root = ../presentation.tex
% =============================================================
%  Subsection — the distributed (wave) representation of the line
%  Part of ac-transmission; \input from ac-transmission.tex,
%  after reactance.tex and before reference-line.tex.
%
%  Fourth step of the chapter's argument: the pi-section is a lumped
%  approximation; the parameters are distributed, so the line obeys a
%  wave equation.  gamma and Z_w, the cosh/sinh solution, and the
%  exact A, B, C, D that replace the nominal-pi entries of
%  two-port.tex.  R', L', G', C' are defined on the chapter's opening
%  frame, line-parameters.tex.  The three frames on gamma/Z_w and the line equations
%  were moved here from power-transmission.tex, where they sat after
%  the load cases.
%  Values from fig/angular_stability_figures.py constants
%  (R' = 56.9 mOhm/km, X' = 0.336 Ohm/km, C' = 12.7 nF/km, 60 Hz).
% =============================================================
\providecommand{\acfig}[1]{ac-transmission/fig/#1}

\intermezzo{Wave Representation}{The line as a distributed system}

% ---------------------------------------------------------------
\begin{frame}{From the $\pi$-section to the wave equation}
  \begin{block}{Why the lumped model runs out}
    {The parameters are distributed along the line, not concentrated\\
    Two quantities describe the distribution completely.}
  \end{block}

  \vspace{0.6em}
  \setlength{\abovedisplayskip}{3pt}\setlength{\belowdisplayskip}{3pt}
  \[
    \underline{\gamma} = \alpha + j\cdot\beta
      = \sqrt{\underline{Z}'\cdot\underline{Y}'}
      = \sqrt{\left(R'+j\cdot\omega\cdot L'\right)\cdot
              \left(G'+j\cdot\omega\cdot C'\right)}
  \]
  \[
    \underline{Z}_w = \sqrt{\frac{\underline{Z}'}{\underline{Y}'}}
      = \sqrt{\frac{R'+j\cdot\omega\cdot L'}{G'+j\cdot\omega\cdot C'}}
  \]

  \vspace{0.5em}
  \centering
  {\color{atGray}\footnotesize Reference line:
  $\alpha = 1.07 \times 10^{-4}$\,Np/km,
  $\beta = 1.27 \times 10^{-3}$\,rad/km,
  $\underline{Z}_w = 266-j\,22\,\Omega$, wavelength 4936\,km.}
\end{frame}

% ---------------------------------------------------------------
\begin{frame}{Voltage and current along the line}
  \begin{columns}[c,onlytextwidth]
    \begin{column}{0.38\textwidth}\centering
      \input{\acfig{two-port-terminals.tex}}
    \end{column}
    \begin{column}{0.58\textwidth}
      {\footnotesize At distance $x$ from the sending end}\\[0.2em]
      {\footnotesize$\begin{aligned}[t]
        \underline{U} &= \underline{U}_R\cdot
                         \cosh\left(\underline{\gamma}\cdot\left(l-x\right)\right)
                       + \underline{Z}_w\cdot\underline{I}_R\cdot
                         \sinh\left(\underline{\gamma}\cdot\left(l-x\right)\right)\\[0.5em]
        \underline{I} &= \frac{\underline{U}_R}{\underline{Z}_w}\cdot
                         \sinh\left(\underline{\gamma}\cdot\left(l-x\right)\right)
                       + \underline{I}_R\cdot
                         \cosh\left(\underline{\gamma}\cdot\left(l-x\right)\right)
       \end{aligned}$}\\[1.2em]
      {\footnotesize At the sending end $\left(x=0\right)$}\\[0.2em]
      {\footnotesize$\begin{aligned}[t]
        \underline{U}_S &= \underline{U}_R\cdot
                           \cosh\left(\underline{\gamma}\cdot l\right)
                         + \underline{Z}_w\cdot\underline{I}_R\cdot
                           \sinh\left(\underline{\gamma}\cdot l\right)\\[0.5em]
        \underline{I}_S &= \frac{\underline{U}_R}{\underline{Z}_w}\cdot
                           \sinh\left(\underline{\gamma}\cdot l\right)
                         + \underline{I}_R\cdot
                           \cosh\left(\underline{\gamma}\cdot l\right)
       \end{aligned}$}
    \end{column}
  \end{columns}
\end{frame}

% ---------------------------------------------------------------
\begin{frame}{The exact two-port parameters}
  \begin{block}{Closing the loop}
    {The same $\underline{A}$, $\underline{B}$, $\underline{C}$,
    $\underline{D}$ as before --- now exact\\
    The nominal $\pi$ was their short-line approximation.}
  \end{block}

  \vspace{0.4em}
  \setlength{\abovedisplayskip}{3pt}\setlength{\belowdisplayskip}{5pt}
  \[
    \underline{A} = \underline{D} = \cosh\left(\underline{\gamma}\cdot l\right),
    \qquad
    \underline{B} = \underline{Z}_w\cdot\sinh\left(\underline{\gamma}\cdot l\right),
    \qquad
    \underline{C} = \frac{\sinh\left(\underline{\gamma}\cdot l\right)}
                         {\underline{Z}_w}
  \]

  \vspace{0.2em}
  \centering
  \renewcommand{\arraystretch}{1.1}
  {\footnotesize
  \begin{tabular}{@{}r r r r r@{}}
    \textbf{$l$} & \multicolumn{1}{c}{\textbf{$\beta\cdot l$}}
      & \multicolumn{1}{c}{\textbf{$\left|B\right|$ nominal $\pi$}}
      & \multicolumn{1}{c}{\textbf{$\left|B\right|$ exact}}
      & \multicolumn{1}{c}{\textbf{error}} \\
    \multicolumn{1}{c}{mi} & \multicolumn{1}{c}{$^\circ$}
      & \multicolumn{1}{c}{$\Omega$} & \multicolumn{1}{c}{$\Omega$}
      & \multicolumn{1}{c}{\%} \\
    \hline
    100 & 11.7 &  54.8 &  54.5 &  $+0.7$ \\
    200 & 23.5 & 109.7 & 106.7 &  $+2.8$ \\
    \atrowhighlight 300 & 35.2 & 164.5 & 154.4 &  $+6.5$ \\
    400 & 46.9 & 219.4 & 195.8 & $+12.0$ \\
    600 & 70.4 & 329.1 & 252.9 & $+30.1$ \\
    \hline
  \end{tabular}}\\[0.4em]
  {\color{atGray}\footnotesize The $\pi$-section holds to roughly 100\,mi;
  beyond that only the wave form is usable.}
\end{frame}
